\chapter{Diagrama de Cableado Explicado (Referencia del Repositorio)}
\label{ch:wiring-diagram-explained}

\section{Diagrama (Incluido)}
\label{sec:wiring-diagram}

% -------------------- FIGURA TAL COMO SE PROPORCIONÓ --------------------
\begin{figure}[H]
    \centering
    \begin{circuitikz}[scale=0.6, transform shape, font=\large]
        \ctikzset{
            label distance=-0pt,
            every label/.style={font=\scriptsize}
        }
        % =====================================================
        % BLOQUE ESP32
        % =====================================================

        \node[draw, minimum width=3cm, minimum height=12.5cm] (esp) at (-3,0) {ESP32};

        % =====================================================
        % LADO IZQUIERDO (ENTRADAS)
        % =====================================================

        % I2C
        \node[left=0.7cm of esp.north west, yshift=-4.8cm] (i2c_sda) {SDA (21)};
        \node[left=0.7cm of esp.north west, yshift=-5.8cm] (i2c_scl) {SCL (22)};

        % Entradas analógicas
        \node[left=2.25cm of esp.center, yshift=5.5cm] (num1) {NUM1 (32)};
        \node[left=2.25cm of esp.center, yshift=4.5cm] (num2) {NUM2 (33)};
        \node[left=2.25cm of esp.center, yshift=3.5cm] (ain1) {AIN1 (34)};
        \node[left=2.25cm of esp.center, yshift=2.5cm] (ain2) {AIN2 (35)};

        % Líneas de conexión con círculos abiertos más grandes
        \draw (num1) -- ++(-2,0) node[ocirc, scale=3.6]{};
        \draw (num2) -- ++(-2,0) node[ocirc, scale=3.6]{};
        \draw (ain1) -- ++(-2.1,0) node[ocirc, scale=3.6]{};
        \draw (ain2) -- ++(-2.1,0) node[ocirc, scale=3.6]{};

        % Indicadores
        \node[left=1cm of esp.south west, yshift=4cm] (ind1) {IND1 (4)};
        \node[left=1cm of esp.south west, yshift=3cm] (ind2) {IND2 (5)};
        \node[left=1cm of esp.south west, yshift=2cm] (ind3) {IND3 (2)};
        \node[left=1cm of esp.south west, yshift=1cm] (ind4) {IND4 (0)};

        % =====================================================
        % LADO DERECHO (SALIDAS)
        % =====================================================

        % Salidas PWM
        \node[right=0.5cm of esp.north east, yshift=-0.8cm] (pwm1) {PWM LX (18)};
        \node[right=0.5cm of esp.north east, yshift=-1.8cm] (pwm2) {PWM LY (19)};
        \node[right=0.5cm of esp.north east, yshift=-2.8cm] (pwm3) {PWM RX (23)};
        \node[right=0.5cm of esp.north east, yshift=-3.8cm] (pwm4) {PWM RY (25)};
        \node[right=0.5cm of esp.north east, yshift=-4.8cm] (pwm5) {PWM KL (26)};
        \node[right=0.5cm of esp.north east, yshift=-5.8cm] (pwm6) {PWM KR (27)};

        % Entradas de switches
        \node[right=2cm of esp.center, yshift=-1.4cm] (sw1) {SW1 (16)};
        \node[right=2cm of esp.center, yshift=-2.4cm] (sw2) {SW2 (17)};
        \node[right=2cm of esp.center, yshift=-3.4cm] (sw3) {SW3 (13)};
        \node[right=2cm of esp.center, yshift=-4.4cm] (sw4) {SW4 (14)};
        \node[right=2cm of esp.center, yshift=-5.4cm] (sw5) {SW5 (12)};
        \node[right=2cm of esp.center, yshift=-6.4cm] (sw6) {SW6 (15)};



        % =====================================================
        % SENSOR I2C
        % =====================================================

        \node[draw, minimum width=1.5cm, minimum height=0.8cm] (sensor) at (-14,2) {I2C};

        \draw (sensor.east) -- ++(0.3,0) |- (i2c_sda);
        \draw (sensor.east) -- ++(0.3,0) |- (i2c_scl);

        % Puntos de inicio de pull-up
        \coordinate (sda_r) at ($(i2c_sda)+(-3,0)$);
        \coordinate (scl_r) at ($(i2c_scl)+(-5,0)$);

        % Resistencias
        \draw (sda_r) to[R,l={\scriptsize $4.7k\Omega$}, label distance=-1pt] ++(0,1.5) coordinate (vbus1);
        \draw (scl_r) to[R,l={\scriptsize $4.7k\Omega$}, label distance=-1pt] ++(0,2.5) coordinate (vbus2);

        % Unir riel Vcc
        \draw (vbus1) -- ++(0,0.5) coordinate (vccbus);
        \draw (vbus2) |- (vccbus);

        % Nodo Vcc
        \draw (vccbus) -- ++(0,0.6) node[VCC33]{\scriptsize 3.3V};

        % =====================================================
        % MONITOR DE VOLTAJE DE BATERÍA (ENTRADA ADC)
        % =====================================================

        % Referencia local para el circuito
        \coordinate (batbase) at ($(ain1)+(9,-14)$);

        % Resistencia superior
        \draw ($(batbase)+(0,1.5)$)
        to[R,l_=\scriptsize $100k$]
        (batbase);

        % Resistencia inferior
        \draw (batbase)
        to[R,l_=\scriptsize $33k$]
        ($(batbase)+(0,-1.5)$);

        % Positivo de la batería
        \draw ($(batbase)+(0,1.5)$) -- ++(0,1) node[battery1]{};

        % Etiqueta VBAT
        \node[left] at ($(batbase)+(0,2.5)$) {\scriptsize VBAT};

        % Negativo de batería a tierra
        \draw ($(batbase)+(0,-1.5)$) -- ++(0,-1) node[bigground]{};

        % Terminal de conexión ADC
        \draw (batbase) -- ++(2,0)
        node[ocirc, scale=3.6, label=right:{AIN2}]{};

        % =====================================================
        % SENSOR DE TEMPERATURA LM35 (AIN2)
        % =====================================================

        % Referencia local para el bloque LM35
        \coordinate (lm35base) at ($(ain2)+(3.5,-12)$);

        % Cuerpo del sensor
        \node[draw, minimum width=1.8cm, minimum height=1.2cm] (lm35) at (lm35base) {LM35};

        % Pin VCC
        \draw (lm35.north) -- ++(0,0.8) node[VCC33]{3.3V};

        % Pin GND
        \draw (lm35.south) -- ++(0,-0.8) node[bigground]{};

        % Pin de salida al terminal del ESP32
        \draw (lm35.east) -- ++(1.5,0)
        node[ocirc, scale=3.6, label=right:{AIN1}]{};

        % =====================================================
        % ENTRADA DE SENSOR LDR (NUM2)
        % =====================================================

        \coordinate (ldrbase) at ($(num2)+(-7,-14)$);

        \begin{scope}[yshift=-7cm]

            % LDR (parte superior)
            \draw ($(ldrbase)+(0,2.5)$)
            to[photoresistor,l_=\normalsize LDR]
            (ldrbase);

            % Resistencia fija
            \draw (ldrbase)
            to[R,l_=\normalsize $10k$]
            ($(ldrbase)+(0,-1.5)$);

            % VCC
            \draw ($(ldrbase)+(0,1.5)$) -- ++(0,1) node[VCC33]{3.3V};

            % GND
            \draw ($(ldrbase)+(0,-1.5)$) -- ++(0,-1) node[bigground]{};

            % Terminal de conexión
            \draw (ldrbase) -- ++(2,0)
            node[ocirc, scale=3.6, label=right:{NUM2}]{};

        \end{scope}
        % =====================================================
        % ENTRADA DE POTENCIÓMETRO (NUM1)
        % =====================================================

        % Referencia local para el bloque del potenciómetro
        \coordinate (potbase) at ($(num1)+(-2,-15.5)$);

        % Potenciómetro vertical
        \draw ($(potbase)+(0,1.5)$)
        to[pR,l_=\scriptsize $10k$, invert]
        ($(potbase)+(0,-1.5)$);

        % Terminal superior a VCC
        \draw ($(potbase)+(0,0.5)$) -- ++(0,1) node[VCC33]{3.3V};

        % Terminal inferior a GND
        \draw ($(potbase)+(0,-0.5)$) -- ++(0,-1) node[bigground]{};

        % Línea horizontal de la flecha del cursor al terminal de conexión
        \draw (potbase) -- ++(2,0)
        node[ocirc, scale=3.6, label=right:{NUM1}]{};



        % =====================================================
        % SALIDAS LED (BUS DE TIERRA COMÚN)
        % =====================================================

        % Bus de tierra
        \draw (5.8,-3) -- (7,-3) node[bigground]{};

        % Circuitos LED
        \draw (sw1) to[R,l^=$330\Omega$, label distance=-3pt] ++(4.5,0)
        to[led] ++(1.2,0) coordinate (led1) -- ++(0,-1);

        \draw (sw2) to[R,l^=$330\Omega$, label distance=-3pt] ++(4.5,0)
        to[led] ++(1.2,0) coordinate (led2) -- ++(0,-1);

        \draw (sw3) to[R,l^=$330\Omega$, label distance=-3pt] ++(4.5,0)
        to[led] ++(1.2,0) coordinate (led3) -- ++(0,-1);

        \draw (sw4) to[R,l^=$330\Omega$, label distance=-3pt] ++(4.5,0)
        to[led] ++(1.2,0) coordinate (led4) -- ++(0,-1);

        \draw (sw5) to[R,l^=$330\Omega$, label distance=-3pt] ++(4.5,0)
        to[led] ++(1.2,0) coordinate (led5) -- ++(0,-1);

        \draw (sw6) to[R,l^=$330\Omega$, label distance=-3pt] ++(4.5,0)
        to[led] ++(1.2,0) coordinate (led6) -- ++(0,0);



        % =====================================================
        % ENTRADAS DE SWITCH (CON PULL-UP Y BUS COMÚN VCC)
        % =====================================================

        % -------- BUS VCC ----------
        \draw (-7.78,-0.29) -- (-14.5,-0.29) node[VCC33]{3.3V};

        % -----------------------------------------------------
        % IND1
        % -----------------------------------------------------
        \coordinate (ind1b) at ($(ind1)+(-1.2,0)$);
        \draw (ind1) -- (ind1b);

        \draw (ind1b) -- ++(0,0.0)
        to[R,l=$10k\Omega$] ++(0,2)
        coordinate (ind1r);

        \coordinate (sw1) at ($(ind1b)+(-1.5,0)$);
        \draw (sw1) node[bigground, line width=0.1pt]{}
        to[push button,l=\small BTN1] (ind1b);


        % -----------------------------------------------------
        % IND2
        % -----------------------------------------------------
        \coordinate (ind2b) at ($(ind2)+(-3.0,0)$);
        \draw (ind2) -- (ind2b);

        \draw (ind2b) -- ++(0,0.7)
        to[R,l=$10k\Omega$] ++(0,2.3)
        coordinate (ind2r);

        \coordinate (sw2) at ($(ind2b)+(-1.5,0)$);
        \draw (sw2) node[bigground, line width=0.1pt]{}
        to[push button,l=\small BTN2] (ind2b);


        % -----------------------------------------------------
        % IND3
        % -----------------------------------------------------
        \coordinate (ind3b) at ($(ind3)+(-4.8,0)$);
        \draw (ind3) -- (ind3b);

        \draw (ind3b) -- ++(0,1.4)
        to[R,l=$10k\Omega$] ++(0,2.55)
        coordinate (ind3r);

        \coordinate (sw3) at ($(ind3b)+(-1.5,0)$);
        \draw (sw3) node[bigground, line width=0.1pt]{}
        to[push button,l=\small BTN3] (ind3b);


        % -----------------------------------------------------
        % IND4
        % -----------------------------------------------------
        \coordinate (ind4b) at ($(ind4)+(-6.6,0)$);
        \draw (ind4) -- (ind4b);

        \draw (ind4b) -- ++(0,2.1)
        to[R,l=$10k\Omega$] ++(0,2.9)
        coordinate (ind4r);

        \coordinate (sw4) at ($(ind4b)+(-1.5,0)$);
        \draw (sw4) node[bigground, line width=0.1pt]{}
        to[push button,l=\small BTN4] (ind4b);
        % =====================================================
        % SERVOS
        % =====================================================
        \node[draw, minimum width=1.4cm, minimum height=1cm, anchor=west]
        (servo1) at ($(pwm1)+(5,0)$) {Servo};

        \node[draw, minimum width=1.4cm, minimum height=1cm, anchor=west]
        (servo3) at ($(pwm3)+(3.6,0)$) {Servo};

        % Señal PWM
        \draw (pwm1) -- (servo1.west);
        \draw (pwm3) -- (servo3.west);

        % Conexiones de alimentación
        \draw (servo1.north) -- ++(0,0.2) node[VCC5]{5V};
        \draw (servo3.north) -- ++(0,0.2) node[VCC5]{5V};

        \draw (servo1.south) -- ++(0,-0.2) node[bigground]{};
        \draw (servo3.south) -- ++(0,-0.2) node[bigground]{};

        % PWM restante
        \draw (pwm5) -- ++(2.5,0) node[ocirc, scale=3.6]{};
        \draw (pwm6) -- ++(2.5,0) node[ocirc, scale=3.6]{};

        \draw (pwm2) -- ++(2.5,0) node[ocirc, scale=3.6]{};
        \draw (pwm4) -- ++(2.5,0) node[ocirc, scale=3.6]{};

            \end{circuitikz}
        \caption{Asignación sugerida de pines del ESP32 (parte 1).}
        \label{fig:full-rc-wiring-a}
    \end{figure}

\clearpage % forces next floats to appear after a page break

       % --- Figure 2: bottom part (your snippet) ---
       \begin{figure}[H]
           \centering
           \begin{circuitikz}[scale=0.6, transform shape, font=\large]
        % =====================================================
        % CONTROL DE MOTOR PWM (PWM KL 26)
        % =====================================================

        % Referencia local
        \coordinate (motorbase) at (-9.5,-22);

        % Símbolo de motor
        \node[draw, circle, minimum size=0.9cm] (motor) at ($(motorbase)+(0,2.0)$) {\small M};

        % MOSFET con etiqueta de parte
        \node[nmos, anchor=D, scale=2.4, transform shape,
        label=right:{\small IRLZ44N}]
        (mos) at ($(motorbase)+(0,-0.5)$) {};

        % Cable desde el drenador del MOSFET a la parte inferior del motor
        \draw (mos.D) -- (motor.south);

        % Cable desde la parte superior del motor a VCC
        \draw (motor.north) -- ++(0,1.8) node[VCC12]{6V to 12V};

        % Fuente del MOSFET a tierra
        \draw (mos.S) -- ++(0,-0.8) node[bigground]{};

        % Diodo flyback
        \draw (mos.D) -- ++(2.0,0)
        to[Do,l=\small 1N5819]
        ($(motor.north)+(2.0,1)$) -- ($(motor.north)+(0,1)$);

        % Resistencia de compuerta
        \draw ($(mos.G)+(-4,0)$)
        to[R,l=\small 220$\Omega$]
        (mos.G);

        % Resistencia pull-down
        \draw ($(mos.G)+(-1,0)$)
        to[R,l=\small 10k]
        ($(mos.G)+(-1,-2.5)$)
        node[bigground]{};

        % Terminal PWM
        \draw ($(mos.G)+(-4,0)$) -- ++(-0.5,0)
        node[ocirc, scale=3.6, label=left:{\small PWM KL (26)}]{};

        % =====================================================
        % CONTROL LED PWM (PWM KR 27)
        % =====================================================

        % Punto de referencia
        \coordinate (pwmkr) at (-14.5,-18);

        % Circuito LED
        \draw (pwmkr)
        to[R,l=\small 330$\Omega$] ++(2,0)
        to[led,l=\small LED] ++(1.5,0)
        -- ++(0,-1.2) node[bigground]{};

        % Línea de control PWM
        \draw (pwmkr) -- ++(-1,0)
        node[ocirc, scale=3, label=left:{\small PWM KR (27)}]{};


        % =====================================================
        % DRIVER DE MOTOR PUENTE H (DRV8833) PARA CONTROL PWM
        % =====================================================

        \node[draw, thick, minimum width=3.8cm, minimum height=4cm] (drv) at (1,-22) {DRV8833};

        % -----------------------------------------------------
        % Motor A
        % -----------------------------------------------------

        \node[draw, circle, minimum size=0.9cm] (motorA) at (5,-23.0) {M};

        \draw (drv.east) ++(1.6,-0.4)
        node[left]{\small AOUT1}
        -| (motorA.north);

        \draw (drv.east) ++(1.6,-1.6)
        node[left]{\small AOUT2}
        -| (motorA.south);

        \node at (6.6,-23.0) {\small DC Motor A};

        % -----------------------------------------------------
        % Motor B
        % -----------------------------------------------------

        \node[draw, circle, minimum size=0.9cm] (motorB) at (5,-20.8) {M};

        \draw (drv.east) ++(1.6,1.8)
        node[left]{\small BOUT1}
        -| (motorB.north);

        \draw (drv.east) ++(1.6,0.6)
        node[left]{\small BOUT2}
        -| (motorB.south);

        \node at (6.6,-20.8) {\small DC Motor B};

        % -----------------------------------------------------
        % Alimentación de motor
        % -----------------------------------------------------

        \draw ($(drv.north)+(0,0.5)$)
        node[below]{\small VM}
        -- ++(0,1)
        node[VCC12]{2.7V to 10.8V};

        % -----------------------------------------------------
        % Alimentación lógica
        % -----------------------------------------------------

        \draw (drv.west) ++(-1.2,1.7)
        node[right]{\small VCC}
        -- ++(-1.5,0)
        node[vcc]{3.3V};

        % -----------------------------------------------------
        % Tierra
        % -----------------------------------------------------

        \draw (drv.south)
        -- ++(0,-0.5)
        node[bigground,label={[yshift=-25]below:{\small GND}}]{};

        % -----------------------------------------------------
        % Control Motor A
        % -----------------------------------------------------

        \draw (drv.west) ++(-1.2,-1)
        node[right]{\small AIN1}
        -- ++(-1.5,0)
        node[ocirc, scale=3.6,label=left:{\small PWM LY (19)}]{};

        \draw (drv.west) ++(-1.2,-1.6)
        node[right]{\small AIN2}
        -- ++(-1.5,0)
        node[ocirc, scale=3.6,label=left:{\small DIR (26)}]{};

        % -----------------------------------------------------
        % Control Motor B
        % -----------------------------------------------------

        \draw (drv.west) ++(-1.2,0.6)
        node[right]{\small BIN1}
        -- ++(-1.5,0)
        node[ocirc, scale=3.6,label=left:{\small PWM RY (23)}]{};

        \draw (drv.west) ++(-1.2,-0.0)
        node[right]{\small BIN2}
        -- ++(-1.5,0)
        node[ocirc, scale=3.6,label=left:{\small DIR2 (27)}]{};

    \end{circuitikz}
    \caption{Asignación sugerida de pines del ESP32 y conexiones de periféricos para la configuración \texttt{FULL\_RC\_MODE}. Los circuitos mostrados son ejemplos y pueden modificarse según los requisitos específicos de la aplicación del usuario.}
    \label{fig:full-rc-wiring}
\end{figure}
% -------------------- FIN FIGURA --------------------

\section{Resumen Rápido}
\label{sec:wiring-quick-summary}

La Figura~\ref{fig:full-rc-wiring} es una referencia práctica de cableado para el proyecto
\textbf{MICSBOL/ESP32\_BT\_Controller}. Muestra:

\begin{itemize}
    \item \textbf{Entradas (lado izquierdo):}
    \begin{itemize}
        \item Bus I2C (SDA/SCL) para sensores o expansión
        \item Entradas analógicas (NUM1/NUM2/AIN1/AIN2) usadas como fuentes de telemetría
        \item Entradas de indicadores digitales (IND1--IND4) usadas para construir una máscara de estado
    \end{itemize}
    \item \textbf{Salidas (lado derecho):}
    \begin{itemize}
        \item Salidas PWM para servos, motores, brillo de LED y zumbador
        \item Salidas de switches (SW1--SW4 en modo directo, SW1--SW6 en modo MCP)
        \item Salidas de pulso/evento usadas para acciones momentáneas (eventos de botón)
    \end{itemize}
\end{itemize}

\noindent En términos de código:
\begin{itemize}
    \item \texttt{control.cpp} mapea campos de \texttt{rcStatePacket} a acciones (dirección, motores, paneo, LED, zumbador, switches).
    \item \texttt{hal\_output.cpp} realiza las escrituras reales de hardware (PWM vía \texttt{ledcWrite}, GPIO vía \texttt{digitalWrite},
    salidas MCP vía \texttt{mcpWriteOutput}, pulsos vía \texttt{pulseStart}).
    \item \texttt{input\_hw.cpp} lee pines ADC e indicadores (usados como fuentes de telemetría).
    \item \texttt{pins.h} define el mapeo de GPIO (y cambia dependiendo del modo de proyecto).
\end{itemize}

\section{Explicación Directa de Cada Grupo (Qué hace)}
\label{sec:wiring-components}

\subsection{I2C (SDA 21, SCL 22)}
Se usa para conectar dispositivos I2C externos. En el repositorio, el mismo bus también se usa en la configuración
\texttt{FULL\_RC\_MODE\_MCP} para comunicarse con un expansor MCP23017 (dirección \texttt{0x20} en \texttt{pins.h}).

\subsection{Entradas analógicas de telemetría (NUM1/NUM2/AIN1/AIN2)}
Estas son entradas ADC leídas en \texttt{input\_hw.cpp} usando \texttt{analogRead(...)}.
Están pensadas para fuentes de telemetría como voltaje de batería (a través de un divisor), sensores de temperatura, sensores de luz,
potenciómetros, etc.

\subsection{Entradas de indicadores digitales (IND1--IND4)}
Estas son entradas digitales (GPIO) que representan estados on/off (fin de carrera, línea de falla, bumper, etc.).
En \texttt{input\_hw.cpp} se combinan en una máscara de bits para que múltiples estados puedan enviarse a la UI Android.

\subsection{Salidas PWM (PWM LX/LY/RX/RY/KL/KR)}
Las salidas PWM son accionadas por \texttt{hal\_output.cpp} usando \texttt{ledcWrite(...)} y se alimentan con valores desde
\texttt{control.cpp}.
Usos típicos en el concepto de cableado del repositorio:
\begin{itemize}
    \item servos (dirección, paneo de cámara)
    \item PWM de velocidad de motor
    \item brillo de LED
    \item intensidad del zumbador
\end{itemize}

\subsection{Salidas de switches (SW1--SW6)}
Estas son salidas digitales controladas por bits de switch recibidos desde la app Android (\texttt{control.cpp} llama a
\texttt{halSetSwitch(index, state)}).
\begin{itemize}
    \item El modo directo proporciona 4 salidas físicas (\texttt{PIN\_SW1}..\texttt{PIN\_SW4}).
    \item El modo MCP proporciona 6 salidas mediante MCP23017 (\texttt{GPB0\_SW1}..\texttt{GPB5\_SW6}).
\end{itemize}

\subsection{Salidas de pulso/evento}
Los eventos de botones Android se manejan por \texttt{controlHandleEvent(eventId)} en \texttt{control.cpp}, lo que dispara
pulsos momentáneos mediante \texttt{halPulseSwitch(index)}. Los pulsos se implementan de forma no bloqueante en \texttt{pulse.cpp}.
Estas salidas son ideales para disparos (bocina, pestillo, reset) más que para estados continuos.

\section{Resumen (Un Párrafo)}
\label{sec:wiring-resume}

El diagrama de cableado agrupa las conexiones del ESP32 exactamente de la manera en que el código del repositorio está organizado:
\texttt{control.cpp} decide qué significa cada entrada, \texttt{hal\_output.cpp} acciona salidas PWM y GPIO,
\texttt{input\_hw.cpp} lee entradas analógicas y de indicadores para telemetría, y \texttt{pins.h} define el mapeo de GPIO.
Al mantener estas responsabilidades separadas, puedes adaptar el hardware (cambiar números de pin, agregar dispositivos I2C, cambiar a MCP23017)
sin reescribir la lógica general de control.